\clearpage

\begin{center}
    {\LARGE\textbf{Kurzfassung}}
\end{center}
\vspace{1em}

Interne Wissens- und Kollaborationsplattformen spielen eine zentrale Rolle in digitalen Organisationsstrukturen, werden jedoch häufig nur eingeschränkt genutzt, wenn \Gls{ia} und \Gls{ux} nicht an den Bedürfnissen der Nutzer ausgerichtet sind. Diese Arbeit untersucht, wie \Gls{ux}-Strategien zur Optimierung der \Gls{ia} beitragen und welche Faktoren die Akzeptanz solcher Systeme beeinflussen.

Die theoretische Grundlage umfasst Konzepte der \Gls{ia}, \Gls{ux}-Methodik, Kognitionspsychologie und organisationalen Veränderungsprozesse. Empirisch basiert die Arbeit auf einer qualitativen Einzelfallstudie mit leitfadengestützten Experteninterviews, ausgewertet mittels qualitativer Inhaltsanalyse.

Die Ergebnisse zeigen, dass \Gls{ux}-Methoden wie Personas, User Journeys und Design Patterns insbesondere dann wirksam sind, wenn sie systematisch, iterativ und organisationsübergreifend implementiert werden. Akzeptanz organisationaler Plattformen ergibt sich demnach aus dem Zusammenspiel von nutzerzentrierter Struktur, konsistenter Gestaltung und begleitender Change-Kommunikation. Die Arbeit leitet daraus praxisorientierte Empfehlungen für die Gestaltung und Weiterentwicklung interner Plattformen ab.
\vfil

\begin{center}
    {\LARGE\textbf{Abstract}}
\end{center}
\vspace{1em}

Internal knowledge and collaboration platforms are increasingly fundamental to digital organizational environments, yet their effectiveness depends on usability, clarity, and user acceptance. This thesis examines how user experience (UX) strategies influence the structure and effectiveness of information architecture in organizational platforms and identifies factors that support sustainable system adoption.

The theoretical foundation integrates concepts from information architecture, \Gls{ux} methodology, cognitive psychology, and organizational change. Empirically, the study applies a qualitative single-case design using semi-structured expert interviews, analyzed through qualitative content analysis.

Results indicate that \Gls{ux} methods – such as personas, user journeys, and design patterns – contribute most effectively when implemented systematically, iteratively, and collaboratively across organizational roles. Platform acceptance emerges as a function of user-centered structure, consistent design principles, and strategic change communication. Based on these findings, the thesis proposes actionable recommendations for designing and evolving user-centered organizational platforms. 